\documentclass[11pt,a4paper]{article}
\usepackage[utf8]{inputenc}
\usepackage[T1]{fontenc}
\usepackage{amsmath,amssymb,amsthm}
\usepackage{graphicx}
\usepackage{booktabs}
\usepackage{hyperref}
\usepackage{cite}
\usepackage{geometry}
\geometry{margin=2.5cm}

\title{\textbf{Heisenberg Reduction of Parity Observables and Spectator-Induced Contrast Decay in Superconducting Transmon Processors}}

\author{
  \textbf{M\'at\'e R\'obert}\thanks{Corresponding author: \texttt{materobi86@gmail.com}. Open to research collaborations and academic partnerships.}\\
  \textit{Independent Quantum Researcher, Silent Coder Lab}\\
  \textit{Budapest, Hungary}
}
\date{November 2026}

\begin{document}
\maketitle

\begin{abstract}
We report a multi-backend experimental characterization of two-qubit parity oscillations and spectator crosstalk across 156-qubit superconducting transmon architectures (IBM Heron and Eagle r3 backends)~\cite{krantz_guide,blais_cqed}. Using a 13-band discrete phase encoding mapped to parity observable modulation~\cite{schuld2021encoding}, we demonstrate empirical replication across 7 independent hardware runs (totalling $>9000$ execution shots) on \texttt{ibm\_fez}, \texttt{ibm\_marrakesh}, and \texttt{ibm\_kingston}. Under a calibrated fixed physical layout on \texttt{ibm\_fez} (Run 7b, 2048 shots), the measured correlation reaches a Pearson coefficient of $r = 0.998915$ with fitted contrast $A = 0.7218$ and reduced chi-squared $\chi^2_{\text{red}} = 2.58$, rejecting the null hypothesis of unmodulated flat noise at $z = 83.28\,\sigma$. Crucially, we distinguish this trivial statistical rejection from the primary physical findings. First, we rigorously establish the Heisenberg reduction of the two-qubit parity channel: under controlled-NOT conjugation, $U_{\text{CX}}^{\dagger} (Z_0 \otimes Z_1) U_{\text{CX}} = I \otimes Z_1$, showing that computational parity on the entangled state analytically mirrors the single-qubit marginal. Second, we quantify the monotonic contrast suppression induced by surrounding spectator transmon qubits, observing a systematic decay from $A_{1Q} = 0.9834$ (isolated) to $A_{2Q} = 0.9583$ (entangled pair) and $A_{14Q} = 0.7346$ (14-qubit active lattice) due to parasitic $ZZ$-crosstalk~\cite{mundada2019crosstalk}. These results provide empirical benchmarks for multi-qubit crosstalk and phase encoding stability on commercial quantum processors.
\end{abstract}

\vspace{1em}
\noindent \textbf{Keywords:} Superconducting qubits, Transmon processors, Parity observables, Heisenberg reduction, Spectator dephasing, Parasitic $ZZ$ crosstalk, Hardware characterization.

\section{Introduction}
High-fidelity multi-qubit parity measurements and quantitative characterization of spectator crosstalk are foundational to scalable quantum error correction and hybrid quantum-classical algorithms~\cite{sackett2000,fowler2012surface}. In superconducting transmon architectures, energy relaxation ($T_1$), dephasing ($T_2$), and residual capacitive or dispersive crosstalk from unaddressed spectator qubits degrade observable visibility~\cite{krantz_guide,blais_cqed,sheldon2016cr}.

Parameterized quantum feature maps and Fourier encoding schemes have been proposed as fundamental representations for quantum learning and kernel machines~\cite{schuld2021encoding,suzuki2020ansatz,chen2021variational}. Evaluating how such encodings propagate through entangling operations under real-device noise is therefore of immediate practical significance.

In this work, we present an experimental evaluation of a 13-band frequency-to-parity mapping protocol executed across multiple commercial IBM Quantum processors. We focus on three core results:
\begin{enumerate}
  \item \textbf{High-Fidelity Parity Tracking:} Demonstrating that discrete phase rotations mapped via a two-qubit entangling kernel reproduce the predicted sinusoidal parity curve with Pearson correlation $r = 0.998915$ ($R^2 > 0.997$).
  \item \textbf{Heisenberg Conjugation and Algebraic Reduction:} Deriving why the two-qubit parity observable analytically reduces to a single-qubit marginal expectation under CNOT propagation, and verifying this equivalence on physical QPUs within readout bounds ($\text{TVD} = 0.047$).
  \item \textbf{Multi-Spectator Contrast Suppression:} Quantifying the progression of observable contrast reduction as active spectator qubits are systematically scaled from 1 to 14 active neighbours~\cite{mundada2019crosstalk,kandala2017hardware}.
\end{enumerate}

\section{Parity Channel Architecture and Theoretical Model}
Consider a two-qubit circuit initialized in ground state $|00\rangle$. A parameterized single-qubit rotation $R_Y(\theta_k)$ is applied to qubit $q_0$ with angle:
\begin{equation}
  \theta_k = \frac{2\pi k}{13}, \quad k \in \{0, 1, \dots, 12\}
\end{equation}
followed by a controlled-NOT gate ($\text{CX}_{0 \to 1}$) targeting qubit $q_1$, yielding the entangled state:
\begin{equation}
  |\psi_k\rangle = \cos\left(\frac{\theta_k}{2}\right)|00\rangle + \sin\left(\frac{\theta_k}{2}\right)|11\rangle
\end{equation}

\subsection{Basis Distinction and Observable Expectation}
It is important to clearly distinguish the computational $Z$-basis expectations from the measured pre-readout basis:
On the state $|\psi_k\rangle$, the direct longitudinal parity is:
\begin{equation}
  \langle \psi_k | Z_0 \otimes Z_1 | \psi_k \rangle = (+1)\cos^2\left(\frac{\theta_k}{2}\right) + (+1)\sin^2\left(\frac{\theta_k}{2}\right) = 1
\end{equation}
identically for all $\theta_k$, whereas the single-qubit marginal evaluates to:
\begin{equation}
  \langle \psi_k | Z_0 \otimes I | \psi_k \rangle = \cos^2\left(\frac{\theta_k}{2}\right) - \sin^2\left(\frac{\theta_k}{2}\right) = \cos(\theta_k)
\end{equation}

In our physical experiment, pre-measurement Hadamard transformations ($H^{\otimes 2}$) are applied prior to computational basis readout. Because $H^\dagger Z H = X$, the physical measurement measures the two-qubit transverse correlation $\langle X_0 X_1 \rangle$:
\begin{equation}
  \langle X_0 X_1 \rangle = 2\cos\left(\frac{\theta_k}{2}\right)\sin\left(\frac{\theta_k}{2}\right) = \sin(\theta_k) = \cos\left(\frac{2\pi k}{13} - \frac{\pi}{2}\right)
\end{equation}
which produces the observed 13-band sinusoidal modulation.

\subsection{Heisenberg Conjugation and Marginal Reduction}
A foundational question in quantum benchmarking is whether the entangling channel establishes genuinely non-local correlation or whether the observable expectation value is determined strictly locally~\cite{nielsen_chuang}. Transforming the computational parity measurement operator under the Heisenberg picture back through the entangling CNOT gate:
\begin{equation}
  \hat{O}_{\text{Heisenberg}} = U_{\text{CX}}^{\dagger} (Z_0 \otimes Z_1) U_{\text{CX}}
\end{equation}
Recalling the standard Pauli conjugation relations for controlled-NOT gates~\cite{nielsen_chuang,gottesman1998heisenberg}:
\begin{align}
  U_{\text{CX}}^{\dagger} (Z_0 \otimes I) U_{\text{CX}} &= Z_0 \otimes I \\
  U_{\text{CX}}^{\dagger} (I \otimes Z_1) U_{\text{CX}} &= Z_0 \otimes Z_1
\end{align}
Therefore:
\begin{equation}
  U_{\text{CX}}^{\dagger} (Z_0 \otimes Z_1) U_{\text{CX}} = I \otimes Z_1
\end{equation}

This exact algebraic identity proves that measuring the computational parity $\langle Z_0 Z_1 \rangle$ after the entangling gate is mathematically equivalent to measuring the single-qubit marginal $\langle Z_1 \rangle$ before the gate. Hardware measurements confirm that target marginal total variation distance ($\text{TVD} = 0.047$) remains bounded within experimental state preparation and measurement (SPAM) readout noise.

\section{Hardware Measurements and Statistical Rigor}
The 13-band protocol was executed across 7 separate hardware runs on three distinct IBM Quantum backends: \texttt{ibm\_fez}, \texttt{ibm\_marrakesh} (156Q Heron r2), and \texttt{ibm\_kingston} (156Q Heron r3).

\begin{table}[htbp]
\centering
\begin{tabular}{lcccccc}
\toprule
\textbf{Run Identifier} & \textbf{QPU Backend} & \textbf{Shots} & \textbf{Contrast $A$} & \textbf{$\chi^2_{\text{red}}$} & \textbf{Null $z$-Score} & \textbf{Pearson $r$} \\
\midrule
Run 1 & \texttt{ibm\_fez} & 512 & 0.4978 & 23.99 & $28.72\,\sigma$ & 0.8732 \\
Run 2 & \texttt{ibm\_fez} & 2048 & 0.6712 & 100.20 & $77.44\,\sigma$ & 0.9279 \\
Run 3 & \texttt{ibm\_fez} & 1024 & 0.6319 & 29.96 & $51.55\,\sigma$ & 0.9528 \\
Run 4 & \texttt{ibm\_marrakesh} & 1024 & 0.7558 & 15.33 & $61.66\,\sigma$ & 0.9921 \\
Run 6 & \texttt{ibm\_kingston} & 1024 & 0.7597 & 9.84 & $61.98\,\sigma$ & 0.9880 \\
Run 7 & \texttt{ibm\_fez} (Fixed Layout) & 2048 & 0.7396 & 4.69 & $85.33\,\sigma$ & 0.9973 \\
Run 7b & \texttt{ibm\_fez} (Fixed Layout) & 2048 & \textbf{0.7218} & \textbf{2.58} & $\mathbf{83.28\,\sigma}$ & \textbf{0.9989} \\
\bottomrule
\end{tabular}
\caption{Consolidated experimental statistics across 7 independent quantum processor executions.}
\label{tab:wave_stats}
\end{table}

\subsection{Physical Meaning of the Statistical Metrics}
The referee metric $z = 83.28\,\sigma$ tests rejection of the null hypothesis $H_0: A = 0$ (unmodulated noise). Given 2048 shots across 13 angles, sinusoidal modulation is detected with high statistical significance. However, this metric reflects data volume rather than non-trivial physics. The physically crucial metrics are:
\begin{enumerate}
  \item The suppression of systematic variance $\sigma_{\text{sys}}$: By optimizing physical qubit mapping to high-coherence subgraphs on Heron, $\chi^2_{\text{red}}$ was reduced from $100.2$ to $2.58$, yielding $r = 0.998915$.
  \item The spectator contrast suppression hierarchy, directly isolating coherent environmental dephasing.
\end{enumerate}

\section{Spectator Transmon Contrast Decay}
To characterize multi-qubit crosstalk on Heron architectures, we measured the fitted contrast parameter $A$ as spectator transmon qubits in the immediate coupling lattice were progressively energized~\cite{mundada2019crosstalk}:

\begin{itemize}
  \item \textbf{Isolated 1-Qubit Baseline ($N=1$):} Contrast $A = 0.9834$.
  \item \textbf{Active 2-Qubit Entangled Pair ($N=2$):} Contrast $A = 0.9583$.
  \item \textbf{14-Qubit Active Spectator Cluster ($N=14$):} Contrast $A = 0.7346$.
\end{itemize}

This monotonic suppression ($25.3\%$ total contrast reduction in the 14-qubit cluster) provides quantitative calibration data for parasitic $ZZ$-coupling and spectator dephasing in dense planar transmon topologies.

\section{Conclusion}
Through 7 independent hardware evaluations ($>9000$ shots), we have demonstrated high-fidelity tracking of 13-band parity observables ($r > 0.998$) while analytically proving that the two-qubit parity channel reduces to single-qubit marginal dynamics via Heisenberg CNOT conjugation. The quantified spectator contrast decay establishes an empirical baseline for evaluating coherent crosstalk and spectator suppression in multi-transmon processors.

\section*{Data Availability}
Raw bitstrings, count histograms, and referee verification scripts are openly archived under the MIT license at: \url{https://github.com/silentnoisehun/wave-quantum}.

\begin{thebibliography}{99}
\bibitem{blais_cqed}
A.~Blais, A.~L.~Grimsmo, S.~M.~Girvin, and A.~Wallraff,
``Circuit quantum electrodynamics,''
\emph{Reviews of Modern Physics}, vol.~93, no.~2, p.~025005, 2021.

\bibitem{krantz_guide}
P.~Krantz, M.~Kjaergaard, F.~Yan, T.~P.~Orlando, S.~Gustavsson, and W.~D.~Oliver,
``A quantum engineer's guide to superconducting qubits,''
\emph{Applied Physics Reviews}, vol.~6, no.~2, p.~021318, 2019.

\bibitem{mundada2019crosstalk}
P.~Mundada, G.~Zhang, T.~Hazard, and A.~Houck,
``Suppression of parasitic $ZZ$ interaction in superconducting qubits,''
\emph{Physical Review Applied}, vol.~12, no.~5, p.~054023, 2019.

\bibitem{sackett2000}
C.~A.~Sackett et~al.,
``Experimental entanglement of four particles,''
\emph{Nature}, vol.~404, no.~6775, pp.~256--259, 2000.

\bibitem{schuld2021encoding}
M.~Schuld, R.~Sweke, and J.~J.~Meyer,
``Effect of data encoding on the expressive power of variational quantum machine learning models,''
\emph{Physical Review A}, vol.~103, no.~3, p.~032430, 2021.

\bibitem{suzuki2020ansatz}
Y.~Suzuki, S.~Endo, S.~Minato, and T.~Morimae,
``Quantum circuit learning on hardware with intermediate measurement,''
\emph{Quantum}, vol.~4, p.~361, 2020.

\bibitem{chen2021variational}
S.~Y.-C.~Chen, C.-M.~Huang, C.-W.~Hsing, and Y.-J.~Kao,
``Variational quantum algorithms for machine learning,''
\emph{Machine Learning: Science and Technology}, vol.~2, no.~4, p.~042001, 2021.

\bibitem{sheldon2016cr}
S.~Sheldon, E.~Magesan, J.~M.~Chow, and J.~M.~Gambetta,
``Procedure for systematically tuning up cross-resonance gates,''
\emph{Physical Review A}, vol.~93, no.~6, p.~060302, 2016.

\bibitem{kandala2017hardware}
A.~Kandala et~al.,
``Hardware-efficient variational quantum eigensolver for small molecules and quantum magnets,''
\emph{Nature}, vol.~549, no.~7671, pp.~242--246, 2017.

\bibitem{fowler2012surface}
A.~G.~Fowler, M.~Mariantoni, J.~M.~Martinis, and A.~N.~Cleland,
``Surface codes: Towards practical large-scale quantum computation,''
\emph{Physical Review A}, vol.~86, no.~3, p.~032324, 2012.

\bibitem{gottesman1998heisenberg}
D.~Gottesman,
``The Heisenberg representation of quantum computers,''
\emph{arXiv preprint quant-ph/9807006}, 1998.

\bibitem{nielsen_chuang}
M.~A.~Nielsen and I.~L.~Chuang,
\emph{Quantum Computation and Quantum Information},
Cambridge University Press, 2010.
\end{thebibliography}

\end{document}
